Una papereria ven llibretes, bolígrafs i carpetes a una escola, sempre al mateix preu. La primera comanda
va ser de 10 llibretes, 20 bolígrafs i 5 carpetes, i va costar 65\,€; la segona, de 20 llibretes, 10
bolígrafs i 10 carpetes, i va costar 85\,€; i la tercera, de 5 llibretes, 30 bolígrafs i 15 carpetes, i
va costar 100\,€.

\begin{apartats}

\apartat[1,25]{0,75}
Planteja un sistema d'equacions que permeti trobar el preu de cada article i escriu-lo en forma
matricial.

\begin{solucio}
Siguin $x$, $y$ i $z$ els preus, en euros, d'una llibreta, un bolígraf i una carpeta. Dividint cada
equació per 5:
\[
\left.\begin{aligned}10x+20y+5z&=65\\20x+10y+10z&=85\\5x+30y+15z&=100\end{aligned}\right\}
\ \Longleftrightarrow\
\begin{pmatrix}2&4&1\\4&2&2\\1&6&3\end{pmatrix}\begin{pmatrix}x\\y\\z\end{pmatrix}=\begin{pmatrix}13\\17\\20\end{pmatrix}.
\]
\end{solucio}

\begin{tria}{informacio}
\itemtria{resol}{1}{1,25}
Resol el sistema pel mètode de Gauss i interpreta'n la solució.

\begin{solucio}
\[
\left(\begin{array}{ccc|c}2&4&1&13\\4&2&2&17\\1&6&3&20\end{array}\right)
\xrightarrow[2F_3-F_1]{F_2-2F_1}
\left(\begin{array}{ccc|c}2&4&1&13\\0&-6&0&-9\\0&8&5&27\end{array}\right)
\xrightarrow{3F_3+4F_2}
\left(\begin{array}{ccc|c}2&4&1&13\\0&-6&0&-9\\0&0&15&45\end{array}\right).
\]
$z=3$; $-6y=-9$, $y=1{,}5$; $2x=13-6-3=4$, $x=2$. Una llibreta val 2\,€, un bolígraf 1,50\,€ i una
carpeta 3\,€.
\end{solucio}

\itemtria{dues-comandes}{1}{1,25}
Si només es coneguessin les dues primeres comandes, es podria saber el preu de cada article? Justifica-ho
amb el teorema de Rouché-Fröbenius, i digues quins preus es poden saber i quins no.

\begin{solucio}
Amb les dues primeres equacions, $\left.\begin{aligned}2x+4y+z&=13\\4x+2y+2z&=17\end{aligned}\right\}$,
$\begin{vmatrix}2&4\\4&2\end{vmatrix}=-12\ne0$: $\operatorname{rang}A=\operatorname{rang}A^*=2<3$. És
compatible indeterminat, i els preus no es poden saber tots.\\
Ara bé, restant a la segona el doble de la primera, $-6y=-9$: el preu del bolígraf sí que es pot saber,
1,50\,€. En canvi, llibreta i carpeta només compleixen $2x+z=7$: per exemple, 2\,€ i 3\,€, o bé 1\,€ i
5\,€, o bé 3\,€ i 1\,€.
\end{solucio}
\end{tria}

\begin{nomesllarg}
\apartat{0,75}
L'any anterior, els tres articles valien un 10\,\% menys, i les tres comandes haurien costat un 10\,\%
menys. Quin sistema en resulta? Quina relació hi ha entre la seva solució i la d'aquest any? Cal
resoldre'l?

\begin{solucio}
La matriu de coeficients és la mateixa i els termes independents es multipliquen per 0,9:
$AX'=0{,}9B$. Si $AX=B$, aleshores $A(0{,}9X)=0{,}9B$: $X'=0{,}9X$ és solució, i és l'única perquè
$|A|=-30\ne0$. No cal resoldre'l: els preus de l'any anterior eren els d'ara multiplicats per 0,9
(1,80\,€, 1,35\,€ i 2,70\,€).
\end{solucio}
\end{nomesllarg}

\end{apartats}
